\subsection{Bilinear and sesquilinear forms. Rank.}

In this section, A denotes a ring (resp. a ring equipped with an antiautomorphism J), E a left A-module, and F a right A-module (resp. a left A-module). One endows A with the structure of an (A, A)-bimodule defined by left multiplications and right multiplications. In this case, a bilinear (resp. right sesquilinear for J) map from E × F into the bimodule A is called a bilinear (resp. right sesquilinear for J) form on E × F.

When E = F (which implies that it is a sesquilinear form), one often says that a sesquilinear form on E × F is a sesquilinear form on E.

Given two left A-modules E and E′, and two sesquilinear forms \(\Phi\) and \(\Phi'\) for J on E and E' respectively, one says that \(\Phi\) and \(\Phi'\) are equivalent if there exists an isomorphism \(u\) of the A-module E onto the A-module E′ such that \(\Phi'(u(x), u(y)) = \Phi(x, y)\) for all \(x\) , \(y\) in E; then \(\Phi\) is the inverse image of \(\Phi'\) with respect to \(u\) and \(u\) , and \(\Phi'\) is the inverse image of \(\Phi\) with respect to \(u^{-1}\) and \(u^{-1}\) (no. 2).

Let \(\Phi\) be a bilinear form on \(E \times F\) (where F denotes a right A-module). The linear maps \(s_{\Phi}\) and \(d_{\Phi}\) associated with \(\Phi\) (no. 1, def. 2) are then maps from E into the dual \(F^{*}\) of F, and from F into the dual \(E^{*}\) of E.

Thus, by definition,

\[
\Phi (x, y) = \left\langle x, d _ {\Phi} (y) \right\rangle = \left\langle y, s _ {\Phi} (x) \right\rangle .\tag{23}
\]

We shall now define the linear maps associated with a sesquilinear form. Let J be an antiautomorphism of A, and \(\Phi\) a sesquilinear form (on the right) for J on E × F (where F denotes a left A-module); set J' = J \(^{-1}\) . The map \(\Phi'\) from F × E into A defined by

\[
\Phi^ {\prime} (y, x) = \Phi (x, y) ^ {J ^ {\prime}}
\]

\[
(x \in \mathrm{E}, y \in \mathrm{F})
\]

is, as is easily seen, a sesquilinear form (on the right) for J' on F × E. By no. 2 (def. 5), the sesquilinear forms Φ and Φ' identify respectively with bilinear forms on E × F' and on F × E'. The maps dΦ and dΦ' associated with the latter are called the maps associated on the right and on the left with the sesquilinear form \(\Phi\) , and are denoted \(d_{\Phi}\) and \(s_{\Phi}\) . We therefore have, by definition:

\[
\Phi (x, y) = \left\langle x, d _ {\Phi} (y) \right\rangle = \left\langle y, s _ {\Phi} (x) \right\rangle^ {J} \quad (x \in \mathrm{E}, y \in \mathrm{F}).\tag{24}
\]

Thus \(d_{\Phi}\) (resp. \(s_{\Phi}\) ) is a linear map from \(\mathbf{F}^{\mathrm{j}}\) into \(\mathbf{E}^{*}\) (resp. of \(\mathbf{E}^{\mathrm{j}'}\) into \(\mathbf{F}^{*}\) ), or equivalently a semilinear map from \(\mathbf{F}\) into \(\mathbf{E}^{*}\) (resp. of \(\mathbf{E}\) into \(\mathbf{F}^{*}\) ) relative to \(\mathbf{J}\) (resp. \(\mathbf{J}'\) ) if one regards \(\mathbf{J}\) (resp. \(\mathbf{J}'\) ) as a ring isomorphism \(\mathbf{A}^{0}\) (the opposite of \(\mathbf{A}\) ) on \(\mathbf{A}\) , and \(\mathbf{F}\) (resp. \(\mathbf{E}\) ) as a \(\mathbf{A}^{0}\) -module on the right.

Formula (24) and def. 6 of no. 3 immediately imply that for every submodule N of F (resp. M of E), we have

\[
\mathbf {N ^ {0}} = \stackrel {- 1} {s _ {\Phi}} (\mathbf {N ^ {\prime}}) \quad (\text { resp. } \mathbf {M ^ {0}} = \stackrel {- 1} {d _ {\Phi}} (\mathbf {M ^ {\prime}}))\tag{25}
\]

where N′ (resp. M′) is the submodule of the dual F* of F (resp. of the dual E* of E) orthogonal to N (resp. M) (Chap. II, § 4, no. 2).

PROPOSITION 4. --- Suppose that A is a field, and let \(\Phi\) be a bilinear form (resp. sesquilinear form with respect to J) on \(E \times F\) ; for \(E/F^{0}\) to be of finite dimension, it is necessary and sufficient that \(F/E^{0}\) be finite-dimensional, and these dimensions are then equal.

Indeed, let \(\Phi_{1}\) be the nondegenerate form associated with \(\Phi\) , on \((\mathrm{E} / \mathrm{F}^{0}) \times (\mathrm{F} / \mathrm{E}^{0})\) (no. 3). Suppose that \(\mathrm{E} / \mathrm{F}^{0}\) is of finite dimension \(n\) ; since the linear map \(d_{\Phi_1}\) of \(\mathrm{F} / \mathrm{E}^{0}\) (resp. \((\mathrm{F} / \mathrm{E}^{0})^{\mathrm{j}})\) into \((\mathrm{E} / \mathrm{F}^{0})^{*}\) is injective, \(\mathrm{F} / \mathrm{E}^{0}\) is of finite dimension \(n' \leqslant n\) ; by considering \(s_{\Phi_1}\) one sees likewise that \(n \leqslant n'\) .

COROLLARY 1. --- Suppose that A is a field and that \(\Phi\) is nondegenerate. For a subspace M of E to be of finite dimension, it is necessary and sufficient that M \(^{0}\) be of finite codimension in F, and one then has codim M \(^{0}\) = dim M, and M \(^{00}\) = M.

As \(F^{0}=\{0\}\) , the first two assertions follow from Prop. 4 applied to the restriction of \(\Phi\) to M × F. Moreover, \(M^{0}\) is the orthogonal of \(M^{00}\) , hence \(M^{00}\) is of finite dimension equal to codim \(M^{0}=dim M\) ; but since \(M^{00}\supset M\) , one has \(M^{00}=M\) .

COROLLARY 2. --- Under the hypotheses of cor. 1, let M, N be two subspaces of E; we then have \((M + N)^{0} = M^{0} \cap N^{0}\) ; if moreover M and N are of finite dimension, we have \((M \cap N)^{0} = M^{0} + N^{0}\) .

The first assertion is trivial. Suppose M and N are finite-dimensional, and let \(G = M^{0} + N^{0}\) ; we have \(G^{0} = M^{00} \cap N^{00} = M \cap N\) by cor. 1; prop. 4 applied to the restriction of \(\Phi\) to \(M \times G\) then shows (since \(M^{0} \subset G\) and \(G^{0} \subset M\) ) that one has dim \(M/(M \cap N) = \dim G/M^{0} = \text{codim } M^{0} - \text{codim } G\) , and since codim \(M^{0} = \dim M\) , we deduce dim \((M \cap N) = \text{codim } G\) . But we also have dim \((M \cap N) = \text{codim } (M \cap N)^{0}\) by cor. 1, and since \(G \subset G^{00} = (M \cap N)^{0}\) we have \(G = (M \cap N)^{0}\) .

Prop. 4 allows us to make the following definition:

DEFINITION 7. --- Let A be a field (resp. a field equipped with an antiautomorphism J), E a left vector space over A, F a right (resp. left) vector space over A, and Φ a bilinear (resp. sesquilinear for J) form on E × F. Suppose that E/F⁰ and F/E⁰ are of finite dimension over A. The rank of Φ is the (finite) common dimension of the vector spaces E/F⁰ and F/E⁰.

When \(E/F^{0}\) and \(F/E^{0}\) are of infinite dimension, one says that \(\Phi\) is of infinite rank.

PROPOSITION 5. --- With the hypotheses and notation as in def. 7, the linear maps \(s_{\Phi}\) and \(d_{\Phi}\) associated with \(\Phi\) have the same rank, and this rank is equal to the rank of the form \(\Phi\) .

Indeed, the kernel of the map \(d_{\Phi}\) of F into E* is evidently E°, hence its rank equals the dimension of F/E°. Likewise, the rank of \(s_{\Phi}\) is equal to the dimension of E/F°.

PROPOSITION 6. --- With the hypotheses and notation of Definition 7, suppose further that E and F have the same finite dimension. Then the following conditions are equivalent:

\begin{enumerate}
\def\labelenumi{\alph{enumi})}
\item
  \(d_{\Phi}\) is injective;
\item
  \(d_{\Phi}\) is surjective;
\item
  \(s_{\Phi}\) is injective;
\end{enumerate}

\(d)\) \(s_{\Phi}\) is surjective;

\begin{enumerate}
\def\labelenumi{\alph{enumi})}
\setcounter{enumi}{4}
\tightlist
\item
  \(\Phi\) is nondegenerate.
\end{enumerate}

Indeed, since E, F, E* and F* have the same finite dimension, a) and b) are equivalent, as are c) and d) (chap. II, § 3, no. 4). Since \(s_{\Phi}\) and \(d_{\Phi}\) have the same rank (prop. 5), a) and c) are equivalent. Since e) is equivalent to the relation \(\mathrm{E}^{0} = \mathrm{F}^{0} = \{0\}\) , it is equivalent to the conjunction of a) and c), whence the equivalence of the stated conditions.

COROLLARY. --- With the hypotheses and notations being those of def. 7, suppose further that E is finite-dimensional and that \(\Phi\) is nondegenerate. Then we have dim E = dim F and, for every basis \((e_i)\) ( \(1 \leqslant i \leqslant \dim E\) ) of E, there exists a basis \((f_i)\) of F such that \(\Phi(e_i, f_k) = \delta_{ik}\) ( \(i, k = 1, \ldots, \dim E\) ).

Indeed, as \(\Phi\) is nondegenerate, we have \(E^{0} = F^{0} = \{0\}\) , whence dim \(E = \dim F\) (prop. 4). It follows (prop. 6) that \(d_{\Phi}\) is an isomorphism of \(F\) (resp. \(F^{j}\) ) onto \(E^{*}\) ; therefore, if \((e_{i}^{*})\) is the dual basis of \((e_{i})\) , the elements \(f_{i} = d_{\Phi}^{-1}(e_{i}^{*})\) form a basis of \(F\) which, in view of formula (23) (resp. formula (24)), satisfies \(\Phi(e_{i}, f_{k}) = \delta_{ik}\) .

It is immediate that, in this corollary, one may interchange the roles of E and F, replacing \(d_{\Phi}\) by \(s_{\Phi}\) in the proof.

Remark. --- Let A be a ring equipped with an antiautomorphism J, and let M and N be right A-modules, and \(\Phi\) a left sesquilinear form for J on M × N (no. 2, Remark); it therefore satisfies the equality

\[
\Phi (x a, y a ^ {\prime}) = a ^ {j} \Phi (x, y) a ^ {\prime} \quad (a, a ^ {\prime} \in A, x \in M, y \in N).
\]

The map \(\Phi'\) from N × M into A defined by \(\Phi'(y, x) = \Phi(x, y)^{J'}\) (where J' = J \(^{-1}\) ) is a left sesquilinear form for J', and \(\Phi\) and \(\Phi'\) are identified with bilinear forms on M \(^{J'}\) × N and N \(^{J'}\) × M, respectively. The maps \(s_{\Phi}\) and \(s_{\Phi'}\) associated with these bilinear forms are called the maps associated on the left and on the right with the sesquilinear form \(\Phi\) , and are denoted \(s_{\Phi}\) and \(d_{\Phi}\) . We therefore have, by definition

\[
\Phi (x, y) = \left\langle y, s _ {\Phi} (x) \right\rangle = \left\langle x, d _ {\Phi} (y) \right\rangle^ {J} \quad (x \in M, y \in N),\tag{26}
\]

and \(s_{\Phi}\) (resp. \(d_{\Phi}\) ) is a linear map from \(\mathbf{M}^{\mathbf{j}}\) into \(\mathbf{N}^{*}\) (resp.

of N \(^{j'}\) into M*). One could easily state and prove the analogues, for the case considered here, of def. 7 and Props. 4, 5, 6.

